\begin{textsample}{NEG Dim 2 – 2010\_c5 – Score -2.00 – 2010\_c5/enade\_2010\_serv\_soc\_e35df1477d.txt}  \label{ex:f2_neg_104}
CAPÍTULO 5 DISTRIBUIÇÃO DOS CONCEITOS 5.1 PANORAMA NACIONAL DA DISTRIBUIÇÃO DOS CONCEITOS A Tabela 5.1 apresenta a quantidade e distribuição de cursos de Serviço Social participantes do ENADE/2010 por faixa de conceito e Grande Região.
Observando -se os dados da Tabela 5.1, nota -se que, dos 297 cursos participantes, 118 cursos ( 39,7 \% ) classificaram -se com conceito 4, o valor modal, que se repetiu nas regiões Sudeste ( 46,0 \% ), Sul ( 45,8 \% e Centro-Oeste ( 53,3 \% ).
As regiões Norte e Nordeste tiveram o conceito 3 como valor modal.
Em âmbito nacional, o conceito 3 foi obtido por 52 cursos, o conceito 5 por 28, o conceito 1 por 13 cursos e, finalmente, o conceito 2 por oito cursos.
Dos 297 cursos, 78 ( 26,3 \% ) não receberam conceito ( SC ).
A região Norte participou com 32 cursos ou 10,8 \% do total nacional.
Nesta região, o conceito 3 foi atribuído a 11 cursos dos 32 participantes, equivalentes a 34,4 \% do total regional.
Um pouco mais de 2/5 dos cursos ( 40,6 \%, correspondendo a 13 cursos ) não receberam conceito.
Somente um curso ( 3,1 \% ) recebeu conceito 5.
Os sete restantes foram avaliados com conceito 4 ( 21,9 \% ).
A região Nordeste participou com 52 cursos, 17,5 \% do total nacional.
Destes, 15 cursos, 28,8 \% em termos regionais, obtiveram conceito 3.
Os conceitos 2 e 4 foram atribuídos a, respectivamente, dois e 12 cursos ( 3,8 \% e 23,1 \% ).
O conceito 1 foi atribuído a um curso ( 1,9 \% ) e o conceito 5 a cinco deles ( 9,6 \% ).
Nessa região, 32,8 \% ou 17 dos cursos ficaram sem conceito.
Na região Sudeste observou -se a predominância do conceito 4 ( 46,0 \% ), correspondente a 64 dos 139 cursos participantes.
O conceito 2 foi atribuído a cinco cursos ( 3,6 \% ) e o conceito 3, a 19 ( 13,7 \% ).
O conceito 1 foi atribuído a nove cursos ( 6,5 \% ) e o conceito 5, também a nove ( 6,5 \% ).
Do total de cursos da região, 33 ( 23,7 \% ) ficaram sem conceito.
Dos 59 cursos participantes da região Sul, 27 ( 45,8 \% ) obtiveram conceito 4, o conceito modal.
O conceito 1 foi atribuído a dois cursos ( 3,4 \% ), o conceito 2 a um deles ( 1,7 \% ), o conceito 3 a cinco cursos ( 8,5 \% ) e o conceito 5, a 12 ( 20,3 \% ).
Os 12 cursos restantes ficaram sem conceito.
Na região Centro-Oeste, foram 15 os cursos participantes, dos quais três ( 20,0 \% ) não receberam conceito.
O conceito modal, nessa região, foi 4 ( com oito cursos ou 53,3 \% ).
Somente um curso ( 6,7 \% ) recebeu conceito 1 e um curso ( 6,7 \% ) recebeu conceito 5, enquanto dois cursos ( 13,3 \% ) obtiveram conceito 3. 5.2 CONCEITOS POR CATEGORIA ADMINISTRATIVA E POR GRANDE REGIÃO A Tabela 5.2 apresenta a distribuição dos cursos participantes do ENADE/2010 de Serviço Social, por Categoria Administrativa, de acordo com os conceitos por eles alcançados, segundo as Grandes Regiões brasileiras.
Dos 297 cursos participantes, 56 ( 18,9 \% ) eram ministrados em instituições públicas e 241 ( 81,1 \% ), em privadas.
De acordo com as informações da Tabela 5.2, as instituições privadas foram as que apresentaram, em termos nacionais, a maior proporção de cursos com conceito 5, 17 cursos em 28 do total dessa categoria.
Dos 241 cursos participantes nesta Categoria Administrativa, o conceito 4 foi atribuído a 107 cursos, o conceito modal.
Receberam conceito 1 quatro cursos e conceito 2, também quatro.
O conceito 3 foi a marca de 45 cursos e o conceito 5 foi obtido por 17 deles, como já comentado.
Ficaram sem conceito 64 cursos.
Na rede pública, o conceito modal foi identificado nos conceitos 4 e 5 ( com 11 cursos em cada ).
Entre os demais cursos participantes, nove receberam conceito 1, quatro o conceito 2 e sete o conceito 3, enquanto 14 cursos ficaram sem conceito nesta categoria.
Na análise por região, observa -se que, na região Norte, as instituições privadas participaram com 24 cursos, dos quais dez ficaram sem conceito.
O conceito modal para as instituições privadas na região foi 3 ( com nove cursos ).
Os cinco outros cursos foram avaliados com conceito 4.
As instituições públicas participaram com oito cursos, dos quais três ficaram sem conceito.
O conceito modal para esta Categoria Administrativa na região Norte ficou entre 3 e 4, ambos com dois cursos.
Ouve, ainda, um curso que obteve conceito 1.
Na região Nordeste, a rede privada concentrou 38 dos 52 cursos participantes, 73,1 \% do total da região.
Nesta Categoria Administrativa na região Nordeste, \textbf{predominaram} os cursos com conceito 5 ( 14 cursos ).
Os demais foram avaliados com os conceitos 2 ( 13 cursos ), 3 ( nove cursos ) e 4 ( dois cursos ).
As instituições públicas, nesta região, participaram com 14 cursos, dos quais um ficou sem conceito.
O conceito modal para esta Categoria Administrativa ficou entre 3, 4 e 5, com três cursos em cada conceito.
Além disso, dois cursos obtiveram conceito 1 e outros dois, conceito 2.
Na região Sudeste, a proporção de cursos da rede privada, 87,1 \%, foi mais elevada do que nas demais regiões brasileiras, correspondendo a 121 dos 139 cursos participantes.
Nesta categoria na região, o conceito modal foi 4 ( 62 cursos ).
Os demais foram avaliados com conceito 1 ( quatro cursos ), conceito 2 ( três cursos ), conceito 3 ( 16 cursos ) e conceito 5 ( oito cursos ).
Na combinação de Categoria Administrativa e Grande Região, 28 cursos não receberam conceito.
Entre os 18 cursos em instituições públicas na região Sudeste, a categoria modal foi o conceito 1 ( com cinco cursos ).
Outros cinco cursos ficaram sem conceito.
Os demais cursos foram avaliados com conceito 2 ( dois cursos ), conceito 3 ( três cursos ), 4 ( dois deles ) e conceito 5 ( um curso ).
As instituições privadas concentraram 46 dos 59 cursos participantes da região Sul, 78,0 \% do total regional.
Desses, 23 foram avaliados no conceito modal, 4.
Os demais foram avaliados com conceito 2 ( um curso ), conceito 3 ( cinco cursos ) e conceito 5 ( sete cursos ).
Na combinação de Categoria Administrativa e Grande Região, dez cursos não receberam conceito.
As instituições públicas nesta região participaram com 13 cursos ( 22,0 \% ), dos quais cinco receberam o conceito modal, 5.
Os demais foram avaliados com conceito 1 ( dois cursos ) e conceito 4 ( quatro cursos ), ao passo que outros dois cursos ficaram sem conceito.
Na região Centro-Oeste, 12 dos 15 cursos participantes eram de instituições privadas ( 80,0 \% em termos regionais ).
Destes, dois não receberam conceito e os demais receberam conceitos 3 e 4 ( com dois e oito cursos ).
Na rede pública, dos três cursos na região, um não recebeu conceito e os outros dois receberam, respectivamente, os conceitos 1 e 5 ( com um curso em cada ). 5.3 CONCEITOS POR ORGANIZAÇÃO ACADÊMICA E POR GRANDE REGIÃO Na Tabela 5.3 encontra -se a distribuição dos conceitos atribuídos aos cursos participantes do ENADE/2010 na Área de Serviço Social, por Organização Acadêmica, segundo as Grandes Regiões brasileiras.
Dos 297 cursos participantes, 126 eram oferecidos em Universidades, 39 em Centros Universitários e 132 em Faculdades.
De acordo com os dados apresentados, a maioria dos cursos avaliados com conceito 5 eram vinculados a Universidades ( 17 dos 28 cursos com este conceito ).
Esse tipo de Organização Acadêmica teve o conceito 4 como valor modal, com 42 cursos.
Além dos 29 cursos que ficaram sem conceito, os demais foram avaliados com conceito 1 ( 11 cursos ), conceito 2 ( seis cursos ), conceito 3 ( 21 cursos ) e conceito 5 ( 17 cursos, como já mencionado ).
Entre os cursos em Centros Universitários, o conceito modal também foi 4 ( com 21 cursos ), e oito cursos ficaram sem conceito.
Os outros cursos neste tipo de Organização Acadêmica receberam os conceitos 1 ( um curso ), 3 ( sete cursos ) e 5 ( dois cursos ).
Nas Faculdades, dos 132 cursos, 41 ficaram sem conceito e 55 receberam o mesmo conceito modal dos demais tipos de Organização Acadêmica, 4.
Os demais receberam conceito 3, atribuído a 24 cursos, conceito 5, referente a nove deles, além dos conceitos 1 e 2, obtidos por um e dois cursos, respectivamente.
Considerando -se separadamente as regiões brasileiras, verifica -se que, na região Norte, as Universidades concentraram nove dos 32 cursos participantes.
Em Universidades dessa região, quatro cursos ficaram sem conceito e os cinco restantes foram avaliados com conceito 3 ( um curso ), conceito 4 ( três cursos – o conceito modal na combinação de Organização Acadêmica e Grande Região ), e conceito 5 ( um curso ).
Os Centros Universitários da região Norte foram representados por seis cursos, sendo que um não recebeu conceito.
Os demais cursos receberam conceito 3 ( três cursos – o conceito modal ) e conceito 4 ( dois cursos ).
As Faculdades participaram com 17 cursos na região Norte.
Destes, oito ficaram sem conceito e os nove demais receberam conceito 3 ( sete cursos – o conceito modal ) e conceito 4 ( dois cursos ).
Na região Nordeste, as Universidades participaram com 23 dos 52 cursos na Área de Serviço Social da região.
Entre os cursos em Universidades, quatro ficaram sem conceito.
O conceito modal nesta categoria foi 3 ( oito cursos ).
Os demais receberam conceito 1 ( um curso ), conceito 2 ( dois cursos ), conceito 4 ( cinco cursos ) e conceito 5 ( três cursos ).
Os Centros Universitários contaram com dois cursos participantes na região, sendo que um não recebeu conceito e o outro obteve o conceito modal, 3.
As Faculdades foram representadas por 27 cursos na região Nordeste, dos quais 12 não receberam conceito.
O conceito modal foi 4 ( com sete cursos ).
Dos cursos restantes, seis receberam conceito 3 e dois receberam conceito 5.
Na região Sudeste as Universidades concentraram 51 dos 139 cursos da região.
Entre os cursos em Universidades na região, o conceito modal foi 4 ( com 16 cursos ) e 15 ficaram sem conceito.
Os demais cursos receberam conceito 1 ( sete cursos ), conceito 2 ( quatro cursos ), conceito 3 ( seis cursos ) e conceito 5 ( três cursos ).
Os Centros Universitários participaram com 25 cursos na região Sudeste, dos quais 15 obtiveram o conceito modal, 4, e cinco ficaram sem conceito.
Os demais receberam conceito 1 ( um curso ), conceito 3 ( dois cursos ) e conceito 5 ( dois cursos ).
As Faculdades foram representadas por 63 cursos na região Sudeste, que se \textbf{distribuíram} nos conceitos 1 ( um curso ), 2 ( um curso ), 3 ( 11 cursos ), 4 ( 33 cursos – o conceito modal ) e 5 ( quatro cursos ), ou ficaram sem conceito ( 13 cursos ).
Dos 59 cursos da região Sul, 35 eram de Universidades, para os quais o conceito modal foi 4 ( com 15 cursos ).
Nesse tipo de organização, quatro cursos ficaram sem conceito e os demais receberam conceito 1 ( dois cursos ), conceito 3 ( cinco cursos ) e conceito 5 ( nove cursos ).
Os Centros Universitários na região Sul tiveram três dos quatro cursos participantes avaliados no conceito 4, o conceito modal, e um curso ficou sem conceito.
Foram 20 os cursos vinculados a Faculdades na região Sul, e destes, sete ficaram sem conceito.
Os restantes receberam conceito 2 ( um curso ) 4 ( nove cursos, o conceito modal ) e conceito 5 ( três cursos ).
Na região Centro-Oeste, oito dos 15 cursos eram de Universidades.
Nesse tipo de organização, dois cursos ficaram sem conceito e o conceito modal foi 4 ( com três cursos ).
Os cursos restantes foram avaliados com conceito 1 ( um curso ), conceito 3 ( um curso ) e conceito 5 ( também um ).
Os Centros Universitários contaram com apenas dois cursos na região Centro-Oeste, que foram avaliados com os conceitos 3 e 4.
Dos cinco cursos em Faculdades da região Centro-Oeste, um ficou sem conceito e quatro receberam conceito 4.

% matched lemmas: distribuir, predominar
\end{textsample}
